% lecture13.tex: Lecture 13: Exoplanets
% Planetary Systems, Kapteyn Institute

\documentclass[aspectratio=169, 11pt]{beamer}

% ── Theme and macros ────────────────────────────────────────
\usepackage{../common/beamerthemeIPS}
% macros.tex: Shared math macros for IPS lecture slides
% Mirrors definitions in book/_config.yml for consistency
% Uses \providecommand to avoid conflicts with unicode-math

% ── Derivatives ─────────────────────────────────────────────
\providecommand{\dv}[2]{\frac{\mathrm{d} #1}{\mathrm{d} #2}}
\providecommand{\dd}{}\renewcommand{\dd}{\mathrm{d}}
\providecommand{\pdv}[2]{\frac{\partial #1}{\partial #2}}

% ── Solar quantities ────────────────────────────────────────
\newcommand{\Msun}{M_{\odot}}
\newcommand{\Rsun}{R_{\odot}}
\newcommand{\Lsun}{L_{\odot}}

% ── Planetary quantities ────────────────────────────────────
\newcommand{\Mearth}{M_{\oplus}}
\newcommand{\Rearth}{R_{\oplus}}
\newcommand{\Mjup}{M_{\mathrm{J}}}
\newcommand{\Rjup}{R_{\mathrm{J}}}

% ── Constants ───────────────────────────────────────────────
\newcommand{\kB}{k_{\mathrm{B}}}

% ── Media links ─────────────────────────────────────────────
% Clickable button that opens the animated or video version of a
% figure, e.g. \mediabutton{https://...gif}{Play animation}
\providecommand{\mediabutton}[2]{\href{#1}{\beamergotobutton{#2}}}

\graphicspath{{figures/}{../common/}{../../book/13_exoplanets/figures/}}

% ── Metadata ────────────────────────────────────────────────
\title[Lecture 13: Exoplanets]{Lecture 13: Exoplanets\\Detection, Demographics, Characterisation}
\subtitle{Planetary Systems}
\author[L13]{Tim Lichtenberg}
\institute{Kapteyn Astronomical Institute\\University of Groningen}
\date{2026}
\titlebgcredit{JWST/NIRSpec transit of WASP-39 b. Credit: Alderson et al.\ (2023)}
\titlebgopacity{0.85}

% ════════════════════════════════════════════════════════════
\begin{document}

% ── Title slide ─────────────────────────────────────────────
\begin{frame}[plain,noframenumbering]
  \titlepage
\end{frame}

% ── Learning objectives ─────────────────────────────────────
\begin{frame}{Learning Objectives}
  By the end of this lecture, you will be able to:
  \vskip8pt
  \begin{itemize}
    \item Describe the main exoplanet detection methods and their biases
    \item Interpret the period-radius diagram and the radius valley
    \item Apply transit and radial-velocity geometry to derive bulk density
    \item Evaluate JWST-era atmospheric results and biosignature claims
  \end{itemize}
  \vskip8pt
  \keyresult{From zero confirmed exoplanets in 1991 to more than 6000 in 2026.}
\end{frame}

% ── Recap ───────────────────────────────────────────────────
\begin{frame}{Recap: Where We Are in the Course}
  \begin{itemize}
    \item Lectures 9--11 surveyed the geology, atmospheres, and interiors of the solar planets
    \item Lecture 12 decoded small bodies and meteorites as planet formation fossils
    \item This lecture turns outward to thousands of exoplanets to test if these patterns are universal
  \end{itemize}
  \vskip8pt
  \textbf{Today:} How do we detect, weigh, and characterise worlds beyond our solar system?
\end{frame}


% ════════════════════════════════════════════════════════════
\sectionimage{hd209458b_first_transit}{HD 209458 b first transit, STARE 1999. Charbonneau et al.\ (2000)}
\section{Historical Context}
% ════════════════════════════════════════════════════════════

% ── Three decades of milestones ─────────────────────────────
\begin{frame}{From Zero to Six Thousand in Three Decades}
  Three decades: from the first detection to a census of over 6000 confirmed worlds.
  \vskip8pt
  \begin{itemize}
    \item \textbf{1992 (PSR B1257+12):} first exoplanets, around a pulsar, by timing (Wolszczan \& Frail)
    \item \textbf{1995 (51 Pegasi b):} first planet around a Sun-like star; a 4.23-day orbit forced migration theory (Mayor \& Queloz)
    \item \textbf{1999 (HD 209458 b):} first transit of an RV planet: a gas-giant radius and bulk density
  \end{itemize}
\end{frame}

% ── HD 209458 b first transit ───────────────────────────────
\begin{frame}{HD 209458 b: First Ground-Based Transit (1999)}
  \begin{columns}[c]
    \column{\recapfigcol}
    \ipsrecapfig{hd209458b_first_transit}

    \column{\recaptextcol}
    \begin{itemize}
      \item STARE photometer, two nights in Sep 1999
      \item Successive transits dropped flux by $\sim 1.7\%$
      \item Independent detection by Henry et al.\ (2000)
      \item Exoplanets became physical objects with measurable radii
    \end{itemize}
  \end{columns}
\end{frame}


% ════════════════════════════════════════════════════════════
\sectionimage{rossiter_mclaughlin}{Doppler signal of a transiting planet (Rossiter-McLaughlin). Triaud (2018)}
\section{Detection Method 1: Radial Velocity}
% ════════════════════════════════════════════════════════════

% ── Radial velocity principle ───────────────────────────────
\begin{frame}{Radial Velocity: Doppler Reflex of the Star}
  \begin{itemize}
    \item Star and planet orbit a common barycentre: $a_\star = (m_p / M_\star)\,a_p$
    \item Periodic Doppler shift of stellar absorption lines yields reflex velocity
    \item Semi-amplitude $K_\star$ (edge-on orbit; set $e = 0$ for a circular one):
  \end{itemize}
  \[
    K_\star = \left(\frac{2\pi G}{P}\right)^{1/3}\!\frac{m_p \sin i}{M_\star^{2/3}}\,\frac{1}{\sqrt{1 - e^2}}
  \]
  \begin{itemize}
    \item Jupiter at 5.2 AU: $K_\star \approx 12.5$ m s$^{-1}$
    \item Earth at 1 AU: $K_\star \approx 9$ cm s$^{-1}$
  \end{itemize}
\end{frame}

% ── Noise floor and inclination degeneracy ──────────────────
\begin{frame}{Noise Floor and the $m \sin i$ Degeneracy}
  RV measures only minimum mass $m_p \sin i$, with precision now limited by stellar jitter rather than spectrographs.
  \vskip6pt
  \begin{itemize}
    \item \textbf{Noise floor:} ESPRESSO reaches 10 cm s$^{-1}$, but stellar granulation and spots produce 0.1--1 m s$^{-1}$ jitter
    \item \textbf{Degeneracy:} Orbit inclination $i$ is unknown, so RV alone yields only a lower limit $m_p \sin i$
    \item \textbf{Breaking degeneracy:} Transits guarantee $\sin i \approx 1$, collapsing the ambiguity to yield the true planet mass
  \end{itemize}
\end{frame}


% ════════════════════════════════════════════════════════════
\sectionimage{transit_geometry}{Transit geometry. Winn (2010)}
\section{Detection Method 2: Transits}
% ════════════════════════════════════════════════════════════

% ── Transit geometry ────────────────────────────────────────
\begin{frame}{Transit Geometry}
  \begin{columns}[c]
    \column{\recapfigcol}
    \ipsrecapfig{transit_geometry}

    \column{\recaptextcol}
    \begin{itemize}
      \item Orbital plane near line of sight $\to$ planet crosses stellar disk
      \item Transit depth: $\delta = (R_p / R_\star)^2$
      \item Geometric probability: $p \approx R_\star / a$
      \item Strong selection bias toward short periods and large $R_p / R_\star$
    \end{itemize}
  \end{columns}
\end{frame}

% ── Quick scales ────────────────────────────────────────────
\begin{frame}{Quick Scales}
  Transit depth scales inversely with the square of stellar radius, making small planets around small stars far easier to detect.
  \vskip8pt
  \begin{itemize}
    \item \textbf{Jupiter around Sun:} $\delta \approx 1.01\%$, easily detected from ground-based telescopes
    \item \textbf{Earth around Sun:} $\delta \approx 84$ ppm, requires ultra-precise space photometry
    \item \textbf{Earth around M dwarf ($0.15\,R_\odot$):} $\delta \approx 4000$ ppm, accessible from ground and space
  \end{itemize}
\end{frame}

% ── Light curve anatomy ─────────────────────────────────────
\begin{frame}{Light Curve Anatomy}
  \begin{columns}[c]
    \column{\recapfigcol}
    \ipsrecapfig{transit_lightcurve_schematic}

    \column{\recaptextcol}
    \begin{itemize}
      \item Four contact times: $t_\mathrm{I}$, $t_\mathrm{II}$, $t_\mathrm{III}$, $t_\mathrm{IV}$
      \item Duration yields impact parameter $b$ and orbital inclination
      \item Ingress and egress shape constrains stellar limb darkening
      \item Transit depth determines planet-to-star radius ratio $R_p / R_\star$
    \end{itemize}
  \end{columns}
\end{frame}

% ── JWST WASP-39 b transit hero ─────────────────────────────
\begin{frame}[plain]
  \begin{center}
    \includegraphics[width=0.95\paperwidth, height=0.78\paperheight, keepaspectratio]{jwst_transit_lightcurve}\\
    \footnotesize JWST/NIRSpec transit photometry of WASP-39 b: transit depth measured to 200--600 ppm per spectroscopic bin, pushing photometric noise below typical exoplanet absorption features. Alderson et al.\ (2023).
  \end{center}
\end{frame}


% ════════════════════════════════════════════════════════════
\sectionimage{chenkipping_mass_radius}{Empirical mass-radius relation. Chen \& Kipping (2017)}
\section{Blackboard: Transit Depth + RV $\to$ Bulk Density}
% ════════════════════════════════════════════════════════════

% ── Setup: combine methods ──────────────────────────────────
\begin{frame}{Why Combine Transit and RV?}
  \begin{itemize}
    \item Transit alone: yields planet radius $R_p$, but no mass
    \item RV alone: yields minimum mass $m_p \sin i$, but no radius
    \item Combined: \textbf{both} for the same planet $\to$ bulk density $\bar\rho_p$
    \item Bulk density turns an astronomical detection into a physical, geological body
  \end{itemize}
  \vskip8pt
  \textit{To the board.}
\end{frame}

% ── Derivation result: density ──────────────────────────────
\begin{frame}{Result: Depth, Semi-Amplitude, Density}
  Transit depth $\delta = (R_p/R_\star)^2$ gives $R_p$; the reflex velocity $K_\star = (2\pi G/P)^{1/3}\, m_p \sin i\, M_\star^{-2/3}$ gives $m_p \sin i$ (board).
  \begin{itemize}
    \item A transit forces $\sin i \approx 1$ and collapses the $m_p \sin i$ degeneracy
    \item Both $R_p$ and $m_p$ measured for the same object
    \item Bulk density follows from elementary geometry:
  \end{itemize}
  \[
    \boxed{\; \bar\rho_p = \frac{m_p}{(4/3)\pi R_p^3} = \frac{3 m_p}{4\pi R_p^3} \;}
  \]
  \vskip4pt
  \keyresult{Transit + RV is the single observational machinery that took exoplanet science from exotic claim to compositional census.}
\end{frame}

% ── Empirical mass-radius diagram ───────────────────────────
\begin{frame}{Density Maps to Composition}
  \begin{columns}[c]
    \column{\recapfigcol}
    \ipsrecapfig{chenkipping_mass_radius}

    \column{\recaptextcol}
    \begin{itemize}
      \item $\bar\rho \approx 5.5\ \mathrm{g\,cm^{-3}}$: rocky terrestrial (Earth, Venus)
      \item $\bar\rho \approx 1.3\ \mathrm{g\,cm^{-3}}$: gas giant (Jupiter)
      \item $\bar\rho \approx 0.5\ \mathrm{g\,cm^{-3}}$: inflated volatile envelope
      \item $\bar\rho \approx 2$--$4\ \mathrm{g\,cm^{-3}}$: sub-Neptunes, unseen in the solar system
    \end{itemize}
  \end{columns}
\end{frame}


% ════════════════════════════════════════════════════════════
\sectionimage{pds70_disk}{PDS 70 disk + 2 protoplanets. Haffert et al.\ (2019)}
\section{Other Detection Methods}
% ════════════════════════════════════════════════════════════

% ── Direct imaging principle ────────────────────────────────
\begin{frame}{Direct Imaging: Spatial Separation of Photons}
  Direct imaging separates planet photons from the stellar glare.
  \vskip8pt
  \begin{itemize}
    \item \textbf{Contrast:} suppress the star by $10^{-9}$ (Jupiter) to $10^{-10}$ (Earth twin) at separation $\theta = a/d$
    \item \textbf{Technology:} extreme adaptive optics plus coronagraphs against diffraction and speckle noise
    \item \textbf{Bias:} young, massive, self-luminous giants on wide orbits ($a > 10$ AU)
  \end{itemize}
\end{frame}

% ── HR 8799 discovery ───────────────────────────────────────
\begin{frame}{HR 8799: Four Imaged Giants}
  \begin{columns}[c]
    \column{\recapfigcol}
    \ipsrecapfig{hr8799_discovery}

    \column{\recaptextcol}
    \begin{itemize}
      \item Keck and Gemini AO discovery of three bound co-moving companions (Marois et al.\ 2008); planet e added in 2010
      \item Four planets confirmed across multiple epochs
      \item Masses $5$--$10\,\Mjup$ at orbital separations $14$--$70$ AU
      \item Young host star (30--60 Myr): the planets still glow from their contraction heat
    \end{itemize}
  \end{columns}
\end{frame}

% ── PDS 70 hero ─────────────────────────────────────────────
\begin{frame}[plain]
  \begin{center}
    \includegraphics[width=\linewidth, height=0.85\paperheight, keepaspectratio]{pds70bc_haffert}\\
    \footnotesize PDS 70 b and c inside the disk gap in three bands: MUSE H$\alpha$ (a), SPHERE K1 (b), NACO L$'$ (c); both protoplanets are actively accreting within the cleared cavity. Haffert et al.\ (2019).
  \end{center}
\end{frame}

% ── Astrometry ──────────────────────────────────────────────
\begin{frame}{Astrometry: Gaia Will Find Jupiter Analogues}
  Astrometry follows the star's reflex wobble on the sky and gives true masses.
  \[
    \alpha = \frac{m_p}{M_\star}\cdot\frac{a_p}{d}
  \]
  \begin{itemize}
    \item \textbf{Signal:} a Jupiter analogue at 10 pc gives $\alpha \sim 500\,\mu$as; an Earth analogue $\sim 0.3\,\mu$as
    \item \textbf{Gaia DR4 (2026):} thousands of cold giants expected; $\sim 10\,\mu$as final precision
    \item \textbf{True mass:} the inclination $i$ is measured, so no $m \sin i$ degeneracy
  \end{itemize}
\end{frame}

% ── Microlensing and timing ─────────────────────────────────
\begin{frame}{Microlensing, Timing, and Detection Biases}
  Each method covers its own part of parameter space, so the archive carries the biases of the methods.
  \vskip6pt
  \begin{itemize}
    \item \textbf{Microlensing (Roman):} cool planets at 1--10 AU and unbound rogue worlds
    \item \textbf{Timing variations (TTVs):} dynamical masses for faint hosts such as TRAPPIST-1
    \item \textbf{Together:} transits and RV probe close-in orbits; imaging, astrometry and microlensing the wide ones
  \end{itemize}
\end{frame}


% ════════════════════════════════════════════════════════════
\sectionimage{petigura_occurrence}{Kepler occurrence rates. Petigura et al.\ (2018)}
\section{Demographics: Kepler and TRAPPIST-1}
% ════════════════════════════════════════════════════════════

% ── Kepler occurrence ───────────────────────────────────────
\begin{frame}{The Kepler Revolution}
  Kepler's computable detection efficiency delivered the first bias-corrected occurrence rates: planets outnumber stars in the Milky Way.
  \vskip8pt
  \begin{itemize}
    \item \textbf{Planets are ubiquitous:} Average of $>1$ planet per main-sequence star (Petigura et al.\ 2018)
    \item \textbf{Small planets dominate:} Super-Earths and sub-Neptunes are common; hot Jupiters are rare ($0.5$--$1\%$)
    \item \textbf{Habitable terrestrial rate:} Occurrence of Earth-sized planets in the habitable zone is $\eta_\oplus \sim 0.4$ (Bryson et al.\ 2021)
  \end{itemize}
\end{frame}

% ── Period-size distribution ────────────────────────────────
\begin{frame}{Period-Size Occurrence Rate}
  \begin{columns}[c]
    \column{\recapfigcol}
    \ipsrecapfig{petigura_occurrence}

    \column{\recaptextcol}
    \begin{itemize}
      \item California-Kepler Survey, after stellar parameter refinement
      \item Small planets common across all orbital periods
      \item Small planets outnumber giants at every period
      \item Fall-off beyond $\sim 100$ days is partly completeness-limited
    \end{itemize}
  \end{columns}
\end{frame}

% ── Eta Earth ───────────────────────────────────────────────
\begin{frame}{$\eta_\oplus$: Earths in Habitable Zones}
  \begin{columns}[c]
    \column{\recapfigcol}
    \ipsrecapfig{bryson_etaearth}

    \column{\recaptextcol}
    \begin{itemize}
      \item Bryson et al.\ (2021) occurrence rate for Sun-like stars
      \item Marginalised over planet radius and stellar instellation
      \item Shaded bands represent $68\%$ and $95\%$ credible intervals
      \item Central estimate $\eta_\oplus \sim 0.4$: terrestrial planets in habitable zones are common
    \end{itemize}
  \end{columns}
\end{frame}

% ── TRAPPIST-1 system ───────────────────────────────────────
\begin{frame}{TRAPPIST-1: Seven Earths in Resonance}
  \begin{columns}[c]
    \column{\recapfigcol}
    \ipsrecapfig{trappist1_transits}

    \column{\recaptextcol}
    \begin{itemize}
      \item M8 ultra-cool dwarf at 12 pc hosting seven Earth-sized worlds within 0.06 AU
      \item Resonant Laplace-like chain is consistent with early inward disk migration (Gillon et al.\ 2017)
      \item Three planets (e, f, g) orbit in the temperate habitable zone (Gillon et al.\ 2017)
    \end{itemize}
  \end{columns}
\end{frame}

% ── TRAPPIST-1 TTVs ─────────────────────────────────────────
\begin{frame}{TRAPPIST-1 Masses from TTVs}
  \begin{columns}[c]
    \column{\recapfigcol}
    \ipsrecapfig{trappist1_ttvs}

    \column{\recaptextcol}
    \begin{itemize}
      \item Multi-planet transit-timing variations reach tens-of-minutes amplitude
      \item Simultaneous N-body dynamical modeling breaks mass degeneracies
      \item Determines individual masses and densities without requiring radial velocities
    \end{itemize}
  \end{columns}
\end{frame}


% ── Break ───────────────────────────────────────────────────
\breakslide


% ════════════════════════════════════════════════════════════
\sectionimage{fulton_period_radius}{Fulton period-radius diagram. Fulton et al.\ (2017)}
\section{Period-Radius Diagram and the Radius Valley}
% ════════════════════════════════════════════════════════════

% ── Fulton period-radius hero ───────────────────────────────
\begin{frame}[plain]
  \begin{center}
    \includegraphics[width=0.95\paperwidth, height=0.78\paperheight, keepaspectratio]{fulton_period_radius}\\
    \footnotesize Period-radius distribution of small Kepler planets after stellar parameter refinement, revealing super-Earth and sub-Neptune populations separated by a deficit at $\sim 1.8\,\Rearth$. Fulton et al.\ (2017).
  \end{center}
\end{frame}

% ── Radius valley bimodality ────────────────────────────────
\begin{frame}{The Radius Valley (Fulton Gap)}
  \begin{columns}[c]
    \column{\recapfigcol}
    \ipsrecapfig{fulton_gap}

    \column{\recaptextcol}
    \begin{itemize}
      \item Statistically significant deficit at $1.5$--$2.0\,\Rearth$
      \item Super-Earth population peaks at $\sim 1.3\,\Rearth$ (bare rocky cores)
      \item Sub-Neptune population peaks at $\sim 2.4\,\Rearth$ (volatile envelopes)
      \item Two stripping mechanisms: photoevaporation (stellar XUV) and core-powered mass loss (interior heat, no XUV needed)
    \end{itemize}
  \end{columns}
\end{frame}

% ── Photoevaporation mechanism ──────────────────────────────
\begin{frame}{Photoevaporation Strips H/He Envelopes}
  Energy-limited hydrodynamic escape driven by high stellar X-ray and EUV irradiation:
  \[
    \dot M \approx \frac{\epsilon\,\pi F_\mathrm{XUV} R_p^3}{G M_p}
  \]
  \begin{itemize}
    \item Young stars saturate XUV flux at $\sim 10^3\times$ modern levels for $\sim 100$ Myr
    \item Strips light primordial H/He envelopes from lower-mass cores at close-in orbits
    \item Owen \& Wu (2013): naturally carves the bimodal radius valley
  \end{itemize}
\end{frame}

% ── Photoevaporation evolution tracks ───────────────────────
\begin{frame}{Photoevaporation Model + Mass-Loss Tracks}
  \begin{columns}[T]
    \column{0.5\textwidth}
    \includegraphics[width=\linewidth, height=0.60\textheight, keepaspectratio]{owen_evaporation_valley}\\
    \footnotesize \textbf{Left}: predicted radius valley
    \column{0.5\textwidth}
    \includegraphics[width=\linewidth, height=0.60\textheight, keepaspectratio]{owen_xuv_massloss}\\
    \footnotesize \textbf{Right}: time evolution of radius and mass; saturated XUV phase drives bimodal outcome (Owen \& Wu 2013).
  \end{columns}
\end{frame}

% ── Valley slope and stripped cores ─────────────────────────
\begin{frame}{Valley Slope and Stripped Cores}
  \begin{columns}[c]
    \column{\recapfigcol}
    \ipsrecapfig{vaneylen_radius_valley}
    \source{Van Eylen et al.\ (2018)}

    \column{\recaptextcol}
    \begin{itemize}
      \item Asteroseismic radii reveal the valley slope shifts to smaller $R_p$ at longer periods
      \item Consistent with photoevaporation and core-powered cooling models
      \item Many close-in super-Earths are stripped cores rather than primordial rocky worlds
    \end{itemize}
  \end{columns}
\end{frame}


% ════════════════════════════════════════════════════════════
\sectionimage{mazeh_neptune_desert}{Hot Neptune desert. Mazeh et al.\ (2016)}
\section{Hot Neptune Desert and Peas-in-a-Pod}
% ════════════════════════════════════════════════════════════

% ── Hot Neptune desert ──────────────────────────────────────
\begin{frame}{The Hot Neptune Desert}
  \begin{columns}[c]
    \column{\recapfigcol}
    \ipsrecapfig{mazeh_neptune_desert}
    \source{Mazeh et al.\ (2016)}

    \column{\recaptextcol}
    \begin{itemize}
      \item Pronounced deficit of short-period planets ($P < 5$ d) in the $10$--$100\,\Mearth$ mass range
      \item \textbf{Bottom edge:} Photoevaporation strips low-mass envelopes into super-Earths
      \item \textbf{Top edge:} Tidal Roche-lobe overflow strips inflated gas giants
    \end{itemize}
  \end{columns}
\end{frame}

% ── Peas-in-a-pod architectures ─────────────────────────────
\begin{frame}{Peas-in-a-Pod Architectures}
  \begin{columns}[T]
    \column{0.5\textwidth}
    \includegraphics[width=\linewidth, height=0.50\textheight, keepaspectratio]{weiss_peas_in_pod}\\
    \footnotesize Adjacent radii: Pearson $r = 0.65$
    \column{0.5\textwidth}
    \includegraphics[width=\linewidth, height=0.50\textheight, keepaspectratio]{weiss_spacing}\\
    \footnotesize Adjacent period ratios: clustered
  \end{columns}
  \vskip4pt
  \begin{itemize}
    \item Planets in compact multi-systems tend to have similar sizes and uniform orbital spacing (Weiss et al.\ 2018)
    \item Favours quiet in-situ growth or migration over chaotic giant impacts
  \end{itemize}
\end{frame}


% ════════════════════════════════════════════════════════════
\sectionimage{obliquity_pathways}{Hot Jupiter migration channels. Albrecht et al.\ (2022)}
\section{Hot Jupiters and Migration}
% ════════════════════════════════════════════════════════════

% ── Migration channels ──────────────────────────────────────
\begin{frame}{Hot Jupiter Migration Channels}
  \begin{columns}[c]
    \column{\recapfigcol}
    \ipsrecapfig[0.6\textheight]{obliquity_pathways}
    \source{Albrecht et al.\ (2022)}

    \column{\recaptextcol}
    \begin{itemize}
      \item \textbf{Disk migration (Type II):} disk torques move the planet inward smoothly, preserving low obliquity (in situ growth at $0.05$ AU fails: too little solid mass, too hot)
      \item \textbf{High-eccentricity (Kozai-Lidov):} distant perturbers pump $e$; tidal circularisation leaves wide misalignments
      \item \textbf{Planet scattering:} instability ejects giants and scrambles alignments
    \end{itemize}
  \end{columns}
\end{frame}

% ── Rossiter-McLaughlin effect ──────────────────────────────
\begin{frame}{Rossiter-McLaughlin Effect}
  \begin{columns}[c]
    \column{\recapfigcol}
    \ipsrecapfig{rossiter_mclaughlin}

    \column{\recaptextcol}
    \begin{itemize}
      \item Transiting planet first occults blueshifted (approaching) hemisphere, then redshifted
      \item Time-resolved line distortion traces sky-projected spin-orbit angle $\lambda$
      \item Diagnostic tool distinguishing disk migration from dynamical scattering (Triaud 2018)
    \end{itemize}
  \end{columns}
\end{frame}

% ── Obliquity distribution ──────────────────────────────────
\begin{frame}{Observed Obliquity Distribution}
  Stellar obliquity distributions show that cool stars realign tidally while hot stars retain widely scattered spin-orbit orientations.
  \vskip6pt
  \begin{columns}[c]
    \column{\recapfigcol}
    \ipsrecapfig{obliquity_distribution}

    \column{\recaptextcol}
    \footnotesize
    \begin{itemize}
      \item \textbf{Cool stars ($T_\mathrm{eff} < 6250$ K):} Tight orbits aligned by tidal realignment in convective envelopes
      \item \textbf{Hot stars ($T_\mathrm{eff} > 6250$ K):} Retain high obliquities due to thin convective zones
      \item \textbf{Primordial scatter:} Wide dispersion at large $a/R_\star$ proves dynamical migration occurs
    \end{itemize}
  \end{columns}
\end{frame}


% ════════════════════════════════════════════════════════════
\sectionimage{dressing_mdwarf_occurrence}{M dwarf occurrence. Dressing \& Charbonneau (2015)}
\section{Compositions and M Dwarfs}
% ════════════════════════════════════════════════════════════

% ── Sub-Neptunes vs super-Earths ────────────────────────────
\begin{frame}{Sub-Neptunes: Water World or H/He?}
  Sub-Neptune bulk densities suffer a severe degeneracy between extended hydrogen-helium envelopes and deep volatile water layers.
  \vskip8pt
  \begin{itemize}
    \item \textbf{Density split:} Super-Earths ($\bar\rho \sim 4$--$8\ \mathrm{g\,cm^{-3}}$) are rocky; sub-Neptunes ($\bar\rho \sim 1$--$3\ \mathrm{g\,cm^{-3}}$) require volatile envelopes
    \item \textbf{Degeneracy:} Density cannot separate a $1\%$ H/He envelope from a $50\%$ water-rich mantle
    \item \textbf{Spectroscopic test:} Transmission spectroscopy is required to determine the atmospheric mean molecular weight
  \end{itemize}
\end{frame}

% ── M dwarf planets and hazards ─────────────────────────────
\begin{frame}{M Dwarfs: Abundant Hosts, Severe Hazards}
  \begin{columns}[c]
    \column{\recapfigcolnarrow}
    \ipsrecapfig{dressing_mdwarf_occurrence}
    \source{Dressing \& Charbonneau (2015)}

    \column{\recaptextcolwide}
    \begin{itemize}
      \item M dwarfs host $\sim 2.5$ small planets each; geometrically favourable for transits and RV
      \item \textbf{XUV desiccation:} Extended pre-main-sequence luminosity can strip entire oceans (Luger \& Barnes 2015)
      \item \textbf{Tidal locking:} 1:1 spin-orbit resonance forces permanent day and night sides
    \end{itemize}
  \end{columns}
\end{frame}


% ════════════════════════════════════════════════════════════
\sectionimage{wasp39b_prism_spectrum}{WASP-39 b PRISM spectrum. Rustamkulov et al.\ (2023)}
\section{JWST Era: Atmospheric Characterisation}
% ════════════════════════════════════════════════════════════

% ── Transmission spectroscopy ───────────────────────────────
\begin{frame}{Transmission Spectroscopy}
  Starlight filtered through an atmosphere during transit encodes molecular absorption.
  \vskip4pt
  Atmospheric scale height:
  \[
    H = \frac{\kB T}{\mu m_u g}
  \]
  Wavelength-dependent transit depth: $\delta(\lambda) \approx (R_p + n_H(\lambda)\, H)^2 / R_\star^2$, so $\Delta\delta/\delta \approx 2 n_H H / R_p$.
  \vskip4pt
  For a hot Jupiter ($T = 1500$ K, $\mu = 2.3$, $g = 25\ \mathrm{m\,s^{-2}}$):
  \begin{itemize}
    \item Scale height $H \approx 200$ km
    \item Absolute spectral modulation: $\sim 200$ ppm (JWST-routine)
    \item Sub-Neptunes: 10--100 ppm; terrestrial atmospheres: $<10$ ppm
    \item High clouds or hazes flatten the spectrum (GJ 1214 b): a flat spectrum does not mean no atmosphere
  \end{itemize}
\end{frame}

% ── WASP-39 b PRISM hero ────────────────────────────────────
\begin{frame}[plain]
  \begin{center}
    \includegraphics[width=0.95\paperwidth, height=0.78\paperheight, keepaspectratio]{wasp39b_prism_spectrum}\\
    \footnotesize JWST/NIRSpec PRISM transmission spectrum of WASP-39 b. Shaded bands are model opacity contributions; data demonstrate detections of $\mathrm{H_2O}$ ($33\sigma$), $\mathrm{CO_2}$ ($28\sigma$), clouds ($21\sigma$), Na ($19\sigma$), and CO ($7\sigma$). Rustamkulov et al.\ (2023).
  \end{center}
\end{frame}

% ── WASP-39 b SO2 photochemistry ────────────────────────────
\begin{frame}{WASP-39 b: First Disequilibrium Photochemistry}
  \begin{columns}[c]
    \column{\recapfigcol}
    \ipsrecapfig{wasp39b_so2_spectrum}

    \column{\recaptextcol}
    \begin{itemize}
      \item Hot Saturn ($0.28\,\Mjup$) early release science target
      \item $4.05\,\mu\mathrm{m}$ absorption feature confirms \textbf{photochemical $\mathrm{SO_2}$}
      \item Requires UV photolysis of $\mathrm{H_2S}$, proving disequilibrium chemistry
      \item First photochemical product identified in an exoplanet atmosphere (Tsai et al.\ 2023)
    \end{itemize}
  \end{columns}
\end{frame}

% ── TRAPPIST-1 b bare rock ──────────────────────────────────
\begin{frame}{TRAPPIST-1 b: Bare Rock Dayside}
  \begin{columns}[T]
    \column{0.5\textwidth}
    \includegraphics[width=\linewidth, height=0.55\textheight, keepaspectratio]{trappist1b_eclipse}\\
    \footnotesize $15\,\mu\mathrm{m}$ secondary eclipse
    \column{0.5\textwidth}
    \includegraphics[width=\linewidth, height=0.55\textheight, keepaspectratio]{trappist1b_emission}\\
    \footnotesize Compared to model atmospheres
  \end{columns}
  \vskip4pt
  $T_B = 503^{+26}_{-27}$ K, consistent with bare rock equilibrium. \textbf{Thick CO$_2$ atmosphere ruled out;} inner M dwarf rocky planets show a widespread lack of thick atmospheres. Greene et al.\ (2023).
\end{frame}

% ── K2-18 b biosignature evaluation ─────────────────────────
\begin{frame}{K2-18 b: Testing a Tentative Biosignature}
  \begin{columns}[c]
    \column{\recapfigcol}
    \ipsrecapfig{k218b_spectrum}
    \source{Madhusudhan et al.\ (2023)}

    \column{\recaptextcol}
    \begin{itemize}
      \item Sub-Neptune in M-dwarf HZ with detected $\mathrm{CH_4}$ and $\mathrm{CO_2}$; disputed tentative DMS ($\lesssim 2\sigma$)
      \item Competing models: Hycean ocean, mini-Neptune envelope, or magma ocean
      \item Exemplifies rigorous real-time vetting required for candidate biosignatures
    \end{itemize}
  \end{columns}
\end{frame}


% ════════════════════════════════════════════════════════════
\sectionimage{kopparapu_hz}{Habitable zone. Kopparapu et al.\ (2013)}
\section{Habitable Zone and Biosignatures}
% ════════════════════════════════════════════════════════════

% ── Classical habitable zone ────────────────────────────────
\begin{frame}{The Classical Habitable Zone}
  \begin{columns}[c]
    \column{\recapfigcol}
    \ipsrecapfig{kopparapu_hz}

    \column{\recaptextcol}
    \begin{itemize}
      \item Inner edge: Simpson-Nakajima runaway greenhouse threshold
      \item Outer edge: Maximum $\mathrm{CO_2}$ greenhouse before $\mathrm{CO_2}$ condenses
      \item Boundaries scale with stellar effective temperature and luminosity
      \item Updated 1D radiative-convective grid from Kopparapu et al.\ (2013)
    \end{itemize}
  \end{columns}
\end{frame}

% ── HZ evolution and false positives ────────────────────────
\begin{frame}{Habitable Zone Evolution and Biosignatures}
  \begin{itemize}
    \item \textbf{Dynamic trajectory:} HZ boundaries shift with stellar evolution; pre-MS M dwarfs subject planets to prolonged runaway desiccation
    \item \textbf{Disequilibrium pairs:} Atmospheric biosignatures demand chemical disequilibrium ($\mathrm{O_2} + \mathrm{CH_4}$), not lone gas species
    \item \textbf{False positives:} Abiotic $\mathrm{O_2}$ can build up via $\mathrm{H_2O}$ photolysis with hydrogen escape or $\mathrm{CO_2}$ photolysis
  \end{itemize}
  \vskip4pt
  \keyresult{Life detection is an inverse problem requiring exclusion of all abiotic pathways.}
\end{frame}


% ════════════════════════════════════════════════════════════
\sectionimage{life_yield}{LIFE detection yield. Quanz et al.\ (2022)}
\section{Comparative Payoff and Frontier Missions}
% ════════════════════════════════════════════════════════════

% ── LIFE / HWO direct imaging hero ──────────────────────────
\begin{frame}
  \centering
  \makebox[\linewidth][c]{\includegraphics[width=0.95\paperwidth, height=0.74\paperheight, keepaspectratio]{life_yield}}\\
  {\footnotesize Predicted detection yields of the LIFE mid-infrared space interferometer for exoplanets within 20 pc, by radius and insolation. Quanz et al.\ (2022).}
\end{frame}


% ════════════════════════════════════════════════════════════
\section*{Summary}
% ════════════════════════════════════════════════════════════

% ── Summary takeaways ───────────────────────────────────────
\begin{frame}{Summary: Today's Course-Level Takeaways}
  \begin{enumerate}
    \item \textbf{Detection biases:} Transit plus radial velocity breaks $m \sin i$ to yield true bulk density
    \item \textbf{Demographics:} Small planets dominate ($\eta_\oplus \sim 0.4$); photoevaporation carves the radius valley at $\sim 1.8\,\Rearth$
    \item \textbf{Architectures:} resonant chains and peas-in-a-pod systems record smooth disk migration and a shared formation environment
    \item \textbf{Characterisation:} JWST detects photochemistry and bare rocks; biosignatures require ruling out false positives
  \end{enumerate}
  \vskip6pt
  \textbf{Next: Lecture 14}
\end{frame}

% ── Final slide ─────────────────────────────────────────────
\begin{frame}[plain,noframenumbering]
  \begin{tikzpicture}[remember picture, overlay]
    \ipscontain{title_background}
    \fill[ipsPrimary, opacity=0.85] (current page.north west) rectangle (current page.south east);
  \end{tikzpicture}
  \vfill
  \begin{center}
    {\color{white}\Huge \textbf{Questions?}}\\[14pt]
    {\color{white}\large Next: \textbf{Lecture 14. Synthesis \& Astrobiology}}\\[10pt]
    {\color{white}\normalsize Lecture notes: \texttt{formingworlds.github.io/IntroductionToPlanetaryScience}}
  \end{center}
  \vfill
\end{frame}

\end{document}
